% TOP_PROBLEM: {"id": 3214, "title": "5-flow conjecture", "category_id": 3, "category": {"id": 3, "name": "graph_theory", "display_name": "Graph Theory", "description": "Problems involving graphs, networks, and their properties.", "slug": "graph-theory", "order_index": 3, "created_at": "2026-07-31T15:26:25.671Z"}, "status": "open", "problem_number": "OPG-126", "set_id": 12}
% FIRST_POSTED: 2026-10-01
% ENTRY_KIND: statement_only
% UnsolvedMath ID 3214; source catalogue code OPG-126
% !TeX program = lualatex
% Definition and sources only; this file is not an LLM solution attempt.
\documentclass[11pt,a4paper]{article}
\usepackage[margin=25mm]{geometry}
\usepackage{fontspec}
\setmainfont{lmroman10-regular.otf}[BoldFont=lmroman10-bold.otf,
 ItalicFont=lmroman10-italic.otf,BoldItalicFont=lmroman10-bolditalic.otf]
\usepackage{amsmath,amssymb,mathtools}
\usepackage{unicode-math}
\setmathfont{latinmodern-math.otf}
\AtBeginDocument{\let\setminus\smallsetminus}
\usepackage{xurl,hyperref}
\hypersetup{colorlinks=true,urlcolor=blue}
\setcounter{secnumdepth}{0}
\setlength{\parindent}{0pt}
\setlength{\parskip}{5pt}
\setlength{\emergencystretch}{3em}
\begin{document}
\section*{5-flow conjecture}\label{3214}
{\small UnsolvedMath ID 3214 · OPG-126 · Graph Theory · OpenGarden}
\subsection{Definitions and mathematical statement}
Let $G=(V,E)$ be a finite undirected
loopless multigraph; parallel edges
are distinguished. A bridge is an
edge whose deletion increases the
number of connected components.
Choose an orientation of every edge.
Write $\delta^+(v)$ and $\delta^-(v)$
for the edges directed out of and
into vertex $v$, respectively.

For an integer $q\geq2$, an integer
$q$-flow is a function $\phi:E\to\mathbb Z$
satisfying conservation at every vertex,
\[
 \sum_{e\in\delta^+(v)}\phi(e)
 =\sum_{e\in\delta^-(v)}\phi(e),
 \qquad |\phi(e)|<q\quad(e\in E).
\]
It is nowhere-zero if additionally
$\phi(e)\neq0$ on every edge.
Reversing the orientation of one
edge and negating its value preserves
these conditions, so existence is
independent of the initially chosen
orientation.

Tutte's $5$-flow conjecture asserts
that every bridgeless $G$ admits a
nowhere-zero integer $5$-flow. Thus
each edge must carry one of
$\pm1,\pm2,\pm3,\pm4$, with exact
integer conservation at every vertex.
Equivalently, an orientation may be
chosen so all values belong to
$\{1,2,3,4\}$.

Bridgelessness is necessary: summing
conservation over one component after
deleting a bridge forces its value
to be zero. Components and isolated
vertices cause no additional condition;
they may be handled independently.
The bound five is uniform over all
orders and degrees, rather than a
bound allowed to depend on the graph.
\subsection{Short English statement}
Can every bridgeless graph carry a conserved nonzero integer flow with each edge value having magnitude at most four?
\subsection{Sources}
\begin{itemize}
\item[S1] ULAM.AI and the UnsolvedMath Contributors, supplied problems.json, record 3214 (OPG-126), OpenGarden collection. Catalogue identity and source transcription; CC BY 4.0.
\url{https://huggingface.co/datasets/ulamai/UnsolvedMath}.
\item[S2] Open Problem Garden, 5-flow conjecture. Original problem and discussion; the mathematical formulation below is expanded from the supplied source record.
\url{http://www.openproblemgarden.org/op/5_flow_conjecture}.
\end{itemize}
\end{document}
